
Al finalizar el desarrollo e integración de \textit{NodeAlert IoT}, se espera obtener un sistema distribuido que supere las capacidades de los prototipos convencionales en términos de robustez, latencia y escalabilidad. Los resultados proyectados se dividen en las siguientes métricas de rendimiento:

\subsection{Eficacia en la Detección y Tiempo de Respuesta}
Se proyecta que el sistema sea capaz de detectar anomalías térmicas y presencia de gases en un tiempo inferior a 500 ms desde la captura de la señal analógica hasta la activación del actuador local (buzzer). El sensor KY-026, al utilizar una interrupción por hardware (flanco de subida en GPIO5), ofrece el menor tiempo de respuesta al detectar flama directamente sin necesidad de pooling \citep{alexander2006}.

\subsection{Determinismo y Estabilidad del RTOS}
Mediante el uso de FreeRTOS, se espera lograr un comportamiento determinístico en los nodos periféricos. La tarea \texttt{taskKY026} (prioridad 4) nunca es desalojada por tareas de menor prioridad como \texttt{taskMqttLoop} (prioridad 1), garantizando la seguridad del sistema incluso bajo alta carga de tráfico de red \citep{valvano2014volume1}.

\subsection{Conectividad y Escalabilidad de la Red}
Se busca validar la arquitectura \textit{Publish/Subscribe} de MQTT, permitiendo la incorporación de hasta 10 nodos periféricos sin degradar significativamente la latencia. Se espera que la pérdida de mensajes sea mínima (inferior al 1\%) mediante el uso de QoS 1 en los tópicos de eventos y comandos \citep{valvano2014volume2}.

\subsection{Interacción y Ubicuidad de la Información}
El sistema debe proporcionar una experiencia de usuario fluida a través del \textbf{Dashboard Web} en React, con visualización de datos históricos y un retraso máximo de 10 segundos (equivalente al intervalo de polling) respecto al tiempo real. Los comandos de override (acknowledge\_alarm, return\_to\_auto) se reflejan en el nodo en menos de 2 segundos gracias a la publicación directa en MQTT.

\subsection{Autonomía y Gestión}
El sistema opera con alimentación por USB/transformador, sin restricciones de consumo energético en esta versión. La optimización se centra en la eficiencia del código: las tareas de FreeRTOS se bloquean cuando no están muestreando, permitiendo que la \texttt{tarea Idle} ponga el núcleo en bajo consumo durante la mayor parte del tiempo de operación \citep{mazidi2014}.

% FOTO REAL: grafico_latencia_alertas.png